\subsection{Front-Loaded Dropout Across Datasets}
\label{sec:extended_benchmarks}

To test how far the scheduling advantage extends beyond CIFAR, we compare MLPs and Transformers on Speech Commands, Jannis, and Tiny ImageNet, together with a financial MLP on FI-2010. These comparisons span audio, tabular classification, and vision, using the depth-12, width-256 configurations detailed in App.~\ref{app:extended_experiments}. Table~\ref{tab:frontloaded_results} reports the test loss and accuracy of validation-selected models, alongside the final-epoch endpoints where they were recorded.

For each run, we retain the checkpoint with the lowest validation cross-entropy. We then choose the reported front-loaded profile by averaging these minimum validation losses across the paired seeds, restricting selection to profiles with complete seed coverage; test loss and accuracy enter only the subsequent comparison with uniform dropout. Most cohorts contain five seeds, while the extended Jannis Transformer comparison contains ten. This is a reanalysis of the saved validation histories, with the original hyperparameter searches and data splits held fixed. The selected profile stays the same across the checkpoint and final columns, so a change of endpoint cannot silently change the model being compared.

The clearest improvement occurs for the Speech Commands MLP: early-step dropout reduces checkpoint test loss by $12.71\%$ relative to uniform, with accuracy increasing from $72.23\%$ to $74.06\%$. Uniform dropout itself reduces loss by $10.82\%$ relative to no dropout, placing this comparison in a regime where regularization already helps. The Transformer gives a smaller, noisier improvement of $4.91\%$ with big-step dropout and an accuracy gain of $1.53$ percentage points; both confidence intervals include zero. For the MLP, the loss-reduction interval is $[11.01,14.36]\%$ and the accuracy-gain interval is $[1.20,2.46]$ percentage points. These audio results use a class-balanced random split of clips from the official training subset, leaving generalization to held-out speakers unresolved.

Jannis provides a tabular comparison with a different learning problem. Validation selects early-step dropout for the MLP, giving a $2.92\%$ reduction in checkpoint test loss and a $1.29$-percentage-point accuracy gain. For the standard Transformer cohort, validation selects big step, whose test-loss improvement is only $0.15\%$. The ten-seed follow-up is effectively a tie: its validation-selected early step increases test loss by $0.27\%$, despite improving accuracy by $0.39$ percentage points. Both Transformer loss intervals include zero, while the MLP interval is $[0.49,5.27]\%$, leaving the MLP as the clearer signal in this tabular comparison.

On Tiny ImageNet, the selected profiles reduce checkpoint loss by $1.94\%$ and $2.17\%$ for the MLP and ViT at 20,000 training examples, and by $1.26\%$ and $1.01\%$ at 80,000 examples. The larger-data MLP still reaches only $7.95\%$ accuracy, compared with $7.09\%$ for uniform, so its relative gain should be read alongside this weak absolute performance. At the final epoch, the corresponding loss reductions are $8.84\%$ and $1.68\%$; the difference from the checkpoint comparison makes clear how much of the final-loss gain develops after the validation optimum. For the 80,000-example ViT, both checkpoint and final loss intervals include zero. The FI-2010 MLP gives a $3.57\%$ checkpoint loss reduction with big-step dropout, with interval $[2.41,4.73]\%$; its $0.67$-percentage-point accuracy gain has an interval spanning zero.

These results compare complete training configurations whose learning rates and dropout strengths follow each cohort's validation protocol; mean dropout can therefore differ between the selected profile and uniform. The fixed-budget CIFAR experiments remain the direct control for allocation. The broader evidence shows a useful but task-dependent advantage, with small effects and losses retained alongside the positive cases. App.~\ref{app:extended_experiments} gives the learning curves and all recorded profiles, including the unrestricted nonuniform comparisons in Table~\ref{tab:benchmark_results}.
